\documentclass[11pt]{article}
\usepackage[utf8]{inputenc}
\usepackage{amsmath,amssymb}
\usepackage[margin=1in]{geometry}
\usepackage{hyperref}
\emergencystretch=3em

\title{Leading-term extraction from the Taylor tree (equal starting exponents)}
\author{Claude, with Daniel Murfet}
\date{2026-09-07}

\begin{document}
\maketitle
\sloppy

\begin{abstract}
Groundwork for Astra~\#14 rank~2 (the unequal-leading-term posterior ratio). For equal starting
exponents $p=(h_1+1)/k_1=(h_2+1)/k_2$ the retained exponent set of the Taylor tree below
$2T\le p+\min(1/k_1,1/k_2)$ is $\{p\}$ (\texttt{polesBelowGen\_eq\_singleton}), so the chart integral
satisfies $Z(N)=N^{-p}(A_p\log N+B_p)+O(N^{-(p+\delta)}(1+\log N))$ with
$\delta=\min(1/k_1,1/k_2)/2$ (\texttt{chart\_leading\_isBigO}). Hence $Z\sim A_pN^{-p}\log N$ when
$A_p\ne0$ (\texttt{chart\_isEquivalent\_log}) and $Z\sim B_pN^{-p}$ when $A_p=0\ne B_p$
(\texttt{chart\_isEquivalent\_const}): the two leading behaviours whose quotient the calculus of
unit~166 handles. Zero \texttt{sorry}/\texttt{axiom}.
\end{abstract}

\section{The things formalised}
{\small
\begin{verbatim}
theorem polesBelowGen_eq_singleton (h₁ h₂ k₁ k₂) (hk₁ hk₂) (p T) (hp₁ hp₂) (hT₁ : p < 2 * T)
    (hT₂ : 2 * T ≤ p + 1 / k₁) (hT₃ : 2 * T ≤ p + 1 / k₂) :
    polesBelowGen h₁ h₂ k₁ k₂ p p T = {p}
def leadingGap (k₁ k₂ : ℕ) : ℝ := min (1 / k₁) (1 / k₂) / 2
theorem chart_leading_isBigO ... :
    (fun N => twoDAmp β b N h₁ h₂ k₁ k₂ (anaAmp (ampCoeff β x y) b)
        - N ^ (-p) * (canonA … p * Real.log N + canonB … p))
      =O[atTop] fun N => N ^ (-(p + leadingGap k₁ k₂)) * (1 + Real.log N)
theorem isLittleO_remainder_log (p δ) (hδ) (A) (hA : A ≠ 0) :
    (fun N => N ^ (-(p + δ)) * (1 + log N)) =o[atTop] powLog A p 1
theorem isLittleO_const_term (p A B) (hA : A ≠ 0) : (fun N => B * N ^ (-p)) =o[atTop] powLog A p 1
theorem isLittleO_remainder_const (p δ) (hδ) (B) (hB : B ≠ 0) :
    (fun N => N ^ (-(p + δ)) * (1 + log N)) =o[atTop] powLog B p 0
theorem chart_isEquivalent_log ... (hA : canonA … p ≠ 0) :
    (fun N => twoDAmp …) ~[atTop] powLog (canonA … p) p 1
theorem chart_isEquivalent_const ... (hA : canonA … p = 0) (hB : canonB … p ≠ 0) :
    (fun N => twoDAmp …) ~[atTop] powLog (canonB … p) p 0
\end{verbatim}
}

\section{Mathematics}
The exponents of $\Lambda(h,k)$ above the common start $p$ are $p+i/k_1$ and $p+j/k_2$ with
$i,j\ge1$, so nothing lies in $(p,p+\min(1/k_1,1/k_2))$; taking $2T=p+\delta$ with $\delta$ half the
gap retains only $p$ and gives the remainder $O(N^{-(p+\delta)}(1+\log N))$ from the general theorem.
The remainder and the constant term are $o(N^{-p}\log N)$, and the remainder is $o(N^{-p})$,
by the elementary limits $N^{-\delta}(1+\log N)/\log N\to0$, $1/\log N\to0$ and
$N^{-\delta}(1+\log N)\to0$.

\section{Paper-facing meaning}
With unit~166 this gives the deterministic leading posterior ratio for equal starting exponents in
all three regimes: $A^\varphi_p\ne0$ (constant limit $A^\varphi_p/A^1_p=y_{\varphi,00}/y_{1,00}$),
$A^\varphi_p=0\ne B^\varphi_p$ (decay like $1/\log N$, the corner-vanishing case Astra~\#14
highlighted), and the power-decay case, which cannot occur at equal starts within this two-term
expansion (it requires the next exponent). The chart-level corollary is the next unit.

\section{Lean notes}
\begin{itemize}
\item \texttt{rintro rfl} on a goal \texttt{α = p} substitutes the theorem variable \texttt{p} away
(``unknown identifier p''); use \texttt{intro hα; rw [hα]}.
\item \texttt{Real.tendsto\_log\_atTop.inv\_tendsto\_atTop} produces the Pi inverse
\texttt{log⁻¹}; rewrite \texttt{Pi.inv\_apply} before \texttt{field\_simp}.
\item \texttt{isLittleO\_iff\_tendsto'} needs the eventual implication ``denominator zero ⇒
numerator zero'', discharged by contradiction since the model term is nonzero for $N>1$.
\item \texttt{Finset.sum\_singleton} after rewriting the retained set to \texttt{\{p\}};
\texttt{IsBigO.trans\_isLittleO} and \texttt{IsLittleO.add} assemble the equivalences via
\texttt{IsLittleO.isEquivalent}.
\end{itemize}

\end{document}
